\chapter{Weights and changes of representatives}
\label{chap:library:weights}

Let $G$ be a finite group and let $\ell$ be a prime. A weight is a
local representation associated with an $\ell$-subgroup of $G$.
The modular formulation uses simple modular
representations that are projective over the normaliser quotient
\cite[p.~369]{Alperin1987}. We use the ordinary character formulation.
See, for example, \cite[p.~90]{Navarro1998}.

Throughout this chapter, $K$ is a field of characteristic zero.
Write $\Irr_K(H)$ for the trace characters of irreducible finite
dimensional $K$-representations of a finite group $H$.
The weight definitions below use these trace characters. They agree
with the usual ordinary character definition of weights when $K$
is a splitting field for the local quotients $N_G(Q)/Q$.

\section{Weights and their local representations}

\begin{definition}
A character $\theta\in\Irr_K(H)$ has \emph{defect zero} if it is
afforded by a simple $K H$-module $V$ such that
\[
 (\dim_K V)_\ell=|H|_\ell.
\]
Denote the set of these characters by
$\operatorname{dz}_\ell(H)$. Since $\theta(1)=\dim_K V$, the numerical
condition is equivalent to $\theta(1)_\ell=|H|_\ell$.

In \emph{character form}, a weight of $G$ is a pair $(Q,\theta)$, where $Q$ is
an $\ell$-radical subgroup of $G$ and
\[
 \theta\in\operatorname{dz}_\ell(N_G(Q)/Q).
\]
In \emph{representation form}, a weight is a pair $(Q,V)$ with the same
condition on $Q$, where $V$ is a simple finite dimensional
$K[N_G(Q)/Q]$-module of defect zero.
\end{definition}

At a fixed subgroup $Q$, two weights in representation form are identified
when their local representations are isomorphic. Two weights in character form
are identified when their local character functions are equal.
The condition that $Q$ is radical is included in these definitions.

\begin{proposition}[Comparison by characters]
Taking characters gives a bijection from weights in representation form up to
local isomorphism to weights in character form. It induces an
$\Aut(G)$-equivariant bijection on $G$-conjugacy classes of weights.
\end{proposition}
\begin{proof}
Trace is invariant under isomorphism, so the map is well defined.
Character separation for finite groups over fields of characteristic
zero gives injectivity, as in Section~\ref{lib:ordinary:realisations}.
A simple representation affording a character of defect zero gives
surjectivity. Transport along a group isomorphism preserves dimension
and simplicity, and taking traces commutes with this transport.
The bijection therefore commutes with automorphisms and descends to
conjugacy classes.
\end{proof}

\begin{implementation}
The local representations use Mathlib's \lean{FDRep}, with
\lean{FDRep.char_iso} giving invariance of the trace under isomorphism.
The Lean comparison maps take
\lean{ModularRep.LocalSplittingCharacterInput} in
\module{ModularRep.WeightCharacterBridge} as an explicit argument.
Its fields require that the trace of every simple local representation
has an irreducible realisation and that equal traces determine simple
local representations up to isomorphism. These properties follow
mathematically from the definition of $\Irr_K$ and the character
separation argument in Section~\ref{lib:ordinary:realisations}.
The record does not require $K$ to be a splitting field.
The character map is
\lean{ModularRep.LocalSplittingCharacterInput.characterOfSimple}.
The two bijections are
\lean{ModularRep.CharacterWeight.isoClassEquiv} and
\lean{ModularRep.CharacterWeight.conjugacyClassEquiv}, for
\lean{ModularRep.CharacterWeight}.
\end{implementation}

\section{Automorphisms and conjugacy classes}

An isomorphism $\phi:H\to H'$ transports a representation by composition
with $\phi^{-1}$. Its character is therefore $\theta\circ \phi^{-1}$.
This is also the convention for transport of class functions in
\cite[Problem~2.1]{Navarro1998}.
In particular, $\alpha\in\Aut(G)$ sends a weight at $Q$ to a weight at
$\alpha(Q)$ through the induced isomorphism
\[
 \bar\alpha:N_G(Q)/Q\longrightarrow
 N_G(\alpha(Q))/\alpha(Q).
\]
The transported local character is $\theta\circ\bar\alpha^{-1}$.
Composition of these transports gives the usual left action of
$\Aut(G)$. The corresponding right action uses $\alpha^{-1}$, so its
subgroup coordinate is $\alpha^{-1}(Q)$.
\begin{implementation}
\lean{ModularRep.OrdinaryIrreducibleCharacter.mapEquiv_apply}
gives the formula $\theta(\phi^{-1}h')$.
Representation transport uses Mathlib's restriction functor
\lean{Action.res} and the equivalence \lean{Action.resEquiv}.
ModularRep constructs the induced normaliser quotient maps and proves
that the resulting weight transports preserve radicality and defect zero.
The actions on weights in representation and character form are defined before
passing to their classes.
\end{implementation}

Conjugation in $G$ gives an equivalence relation on weights.
Write $\Alp(G)$ for the set of weights in character form modulo this relation.
Automorphisms act on $\Alp(G)$ because they preserve conjugacy.
Local isomorphism and conjugacy in $G$ have different roles:
local isomorphism keeps $Q$ fixed, whereas conjugation may replace it
by another subgroup.

\begin{proposition}[Weights with a fixed radical subgroup]
\label{lib:weights:representative}
Fix an $\ell$-radical subgroup $Q\leq G$. Let $\Alp(G)_{[Q]}$ be the
set of $G$-conjugacy classes of weights in character form whose radical
subgroup is conjugate to $Q$. Then
\[
 \operatorname{dz}_\ell(N_G(Q)/Q)\longrightarrow\Alp(G)_{[Q]},
 \qquad \theta\longmapsto[(Q,\theta)]
\]
is a bijection.
\end{proposition}
\begin{proof}
Every class in $\Alp(G)_{[Q]}$ has a representative at $Q$, obtained
by conjugating its radical subgroup to $Q$. This proves surjectivity.

If $(Q,\theta)$ and $(Q,\eta)$ are conjugate by $g\in G$, then
$g\in N_G(Q)$. The induced automorphism of $N_G(Q)/Q$ is conjugation
by $gQ$. Characters are constant on conjugacy classes, so it fixes
$\theta$. Hence $\eta=\theta$, proving injectivity. The same argument
shows that the character obtained by transporting to $Q$ is
independent of the chosen conjugating element.
\end{proof}
\begin{implementation}
The bijection is
\lean{ModularRep.CharacterWeight.localDefectZeroEquivWeightRadicalFibre}.
Its restriction to a block is
\lean{representativeDZEquivWeightBlockRadicalFibre} in the same
namespace, with the local block induction hypotheses described below.
\end{implementation}

\section{The block of a weight}

Fix compatible ordinary and modular character conventions for the
local groups, including the splitting and root identifications needed
for their block assignments.
Let $(Q,\theta)$ be a weight. Inflate $\theta$ to $N_G(Q)$ and let
$b$ be its block. The weight belongs to an ambient block $B$ when
Brauer block induction is defined and $b^G=B$
\cite[p.~90]{Navarro1998}.

Recall the definition of this induction. If $H\leq G$ and
$\lambda_b:Z(kH)\to k$ is the central character of a block $b$,
extend it linearly to $Z(kG)$ by
\[
 \lambda_b^G(C^+)=\lambda_b((C\cap H)^+)
\]
for each conjugacy class $C$ of $G$. Here $C^+$ denotes the sum of
the elements of $C$ in the group algebra, and $k$ is algebraically
closed of characteristic $\ell$.
If $\lambda_b^G$ is an algebra homomorphism, it is the central
character of a unique block, denoted by $b^G$.
This is Brauer's definition of induced block
\cite[p.~87]{Navarro1998}.
For an $\ell$-subgroup $P$, the inclusions
\[
 P C_G(P)\leq H\leq N_G(P)
\]
ensure that induction is defined for every block of $H$
\cite[Theorem~4.14]{Navarro1998}.

Write $\operatorname{Bl}_\ell(G)$ for the set of $\ell$-blocks of $G$.
Suppose the local block assignment, inflation and block induction
commute with automorphisms. Suppose also that inner automorphisms fix
ambient blocks. Then the ambient block assigned to a weight depends
only on its $G$-conjugacy class, and the resulting map
$\Alp(G)\to\operatorname{Bl}_\ell(G)$ is equivariant.
The corresponding compatibility of Brauer block induction with group
isomorphisms is stated in \cite[Problem~4.4]{Navarro1998}.

\begin{implementation}
\lean{ModularRep.CharacterWeight.LocalBlockInductionOperations} in
\module{ModularRep.CharacterWeightBlockAssignment} supplies the local
character block and inflation. It also supplies finite complete block
decompositions of $G$ and the relevant normalisers, together with their
block central characters, and multiplicativity of the induced central
function. These data determine the ambient block.
\lean{ModularRep.CharacterWeight.LocalBlockInductionSource} adds
compatibility with automorphisms and invariance under inner
automorphisms. Its maps
\lean{LocalBlockInductionSource.weightBlock} and
\lean{LocalBlockInductionSource.weightBlock_transport}
give the assignment on conjugacy classes and its equivariance.
The identification of the supplied operations with ordinary and
modular block theory is a hypothesis.
\end{implementation}

\begin{example}[Local maps from a global correspondence]
Let $\Omega:X\to\Alp(G)$ be a bijection. For each representative $Q$
of a radical conjugacy class, set
\[
 X_Q=\{x\in X:\Omega(x)\in\Alp(G)_{[Q]}\}.
\]
These subsets partition $X$. Restricting $\Omega$ and applying
Proposition~\ref{lib:weights:representative} gives a bijection
\[
 X_Q\longrightarrow\operatorname{dz}_\ell(N_G(Q)/Q).
\]
If $\Omega$ is equivariant, these local maps commute with transport
of subgroups and characters.
The construction for Brauer characters with a fixed compatible root
embedding is in
\module{ModularRep.PaperProofs.IntrinsicGlobalToRepresentativeMaps}.
\end{example}

\newpage
\section{Quotient by the \texorpdfstring{$\ell$}{ell}-core}

\begin{proposition}[Transport through the normal $\ell$-core]
Let $P=O_\ell(G)$. A weight $(Q,V)$ of $G$ in representation form determines
a weight in representation form of $G/P$ with radical subgroup $Q/P$ and
local representation transported along
\[
 N_G(Q)/Q\simeq N_{G/P}(Q/P)/(Q/P).
\]
Conversely, a weight in representation form of $G/P$ determines one of $G$
by taking the full inverse image of its radical subgroup and
transporting the local representation along the inverse isomorphism.
Both constructions preserve the dimension of the local representation.
\end{proposition}
\begin{proof}
By Proposition~\ref{lib:groups:core-contained}, every radical subgroup
$Q$ contains $P$. Proposition~\ref{lib:groups:radical-quotient}
identifies radicality before and after quotienting by $P$.
The normaliser quotient isomorphism then transports the local
representation. An isomorphism of groups preserves simplicity,
dimension and group order, and hence preserves defect zero.
The same reasoning applies to the full inverse image of a radical
subgroup of $G/P$.
\end{proof}
\begin{implementation}
The two constructions are
\lean{ModularRep.RepresentationWeight.quotientPCore} and
\lean{ModularRep.RepresentationWeight.liftFromQuotientPCore}
in \module{ModularRep.NormalCoreTransport}.
They construct individual weights. An inverse equivalence on
conjugacy classes is a separate assertion.
\end{implementation}

A comparison for weights in a specified block also requires
compatibility of the character conventions and block assignments on
the two local quotients. For applications involving modular character
triples, the required triple relation must hold for the transported
local character, with the specified roots and block operations.
